\section{定積分の定義}\label{節:リーマン和}
\chapterref{章:導入}では,\,$y=x^2$ の下の面積を求めるために区間を等しい幅に分け,左端の高さで作った和 $L_n$ と右端の高さで作った和 $U_n$ で面積を挟んだ.\,$x^2$ は増加するので,小区間のどの点の高さを使っても,長方形の和は $L_n$ と $U_n$ の間に入る.では,増加するとは限らない関数や,幅の等しくない分け方ではどうだろうか.同じ図形の面積が,幅の分け方や高さを選ぶ点によって変わっては困る.そこで,分け方と点の選び方を明示し,それらに依存しない極限が存在するときに定積分を定義しよう.

$a<b$ とし,区間 $[a,b]$ を
\[
 a=x_0<x_1<\cdots<x_n=b
\]
と分ける.この点の組を\textbf{分割} $P$ といい,第 $k$ 小区間の幅を $\Delta x_k=x_k-x_{k-1}$ と書く.さらに,各小区間 $[x_{k-1},x_k]$ の中から点 $\xi_k$ を1つずつ選ぶ.\,$[a,b]$ で定義された関数 $f$ について,
\[
 S(P,\xi)=\sum_{k=1}^{n}f(\xi_k)\Delta x_k
\]
を作れる.これを\textbf{リーマン和}という.\,$f(\xi_k)$ が正なら高さと幅の積であり,負なら符号を付けた寄与として足し合わせる.

\begin{figure}[htbp]
\centering
\begin{tikzpicture}[x=3cm,y=1.8cm,>=Stealth]
 \filldraw[fill=blue!12,draw=blue!55] (0,0) rectangle (0.8,1.1225);
 \filldraw[fill=blue!12,draw=blue!55] (0.8,0) rectangle (1.8,1.01);
 \filldraw[fill=blue!12,draw=blue!55] (1.8,0) rectangle (3,1.5625);
 \draw[->] (-0.1,0)--(3.35,0) node[right]{$x$};
 \draw[->] (0,-0.13)--(0,2.4) node[above]{$y$};
 \draw[very thick,domain=0:3.1,samples=80] plot (\x,{1+0.25*(\x-1)^2});
 \foreach \xx/\hh/\name in {0.3/1.1225/1,1.2/1.01/2,2.5/1.5625/3}{
  \draw[blue,dashed] (\xx,0)--(\xx,\hh);
  \fill[blue] (\xx,\hh) circle(1.7pt);
  \node[blue] at (\xx,-0.45) {$\xi_{\name}$};
 }
 \node[below left] at (0,0) {$a=x_0$};
 \node[below] at (0.8,0) {$x_1$};
 \node[below] at (1.8,0) {$x_2$};
 \node[below right] at (3,0) {$x_3=b$};
 \node at (2.7,2.28) {$y=f(x)$};
\end{tikzpicture}
\caption{幅が異なっていてもよく,各小区間のどの点の高さを使ってもよい.}
\label{図:リーマン和}
\end{figure}

図\ref{図:リーマン和}のように,分割を等間隔にする必要はなく,代表点 $\xi_k$ を端点にする必要もない.分割を細かくするとは,単に小区間の個数を増やすことではない.1つだけ幅の広い区間が残っている可能性を除くため,最大の幅
\[
 \|P\|=\max_{1\le k\le n}\Delta x_k
\]
を $0$ に近づける.

\begin{definition}{定積分}
$[a,b]$ 上の有界な関数 $f$ に対し,\,$\|P\|\to0$ のとき,分割と代表点の選び方によらずリーマン和が同じ有限値 $I$ に近づくなら,\,$f$ は $[a,b]$ で\textbf{積分可能}であるという.その値を\textbf{定積分}と呼び,
\[
 \int_a^b f(x)\dd x=I
\]
と書く.
\end{definition}

正確には,任意の $\varepsilon>0$ に対して,ある $\delta>0$ を選び,\,$\|P\|<\delta$ を満たすすべての分割とすべての代表点について $|S(P,\xi)-I|<\varepsilon$ とできることである.特定の分け方で極限が求まるだけでは,この定義を満たすことまでは証明できない.

記号 $\int$ はリーマン和の $\sum$ に,\,$\dd x$ は小区間の幅 $\Delta x_k$ に対応する.\,$\dd x$ の $x$ は,どの変数に沿って幅を足しているかを表している.\,$\int_a^b f(t)\dd t$ と文字を替えても同じ数である.\,$x$ や $t$ はこの積分の内部で用いる変数であり,\,$a,b$ が固定されていれば,定積分全体は1つの数になる.
